% distinction_reorganization.tex
% XeLaTeX
\documentclass[11pt, a4paper]{article}

\usepackage{fontspec}
\setmainfont{TeX Gyre Pagella}

\usepackage{microtype}
\usepackage{geometry}
\geometry{
  a4paper,
  top=2.8cm, bottom=3cm,
  left=3.2cm, right=3.2cm
}

\usepackage{setspace}
\setstretch{1.18}

\usepackage{abstract}
\renewcommand{\abstractnamefont}{\normalfont\bfseries}
\renewcommand{\abstracttextfont}{\normalfont}

\usepackage{hyperref}
\hypersetup{
  colorlinks=true,
  linkcolor=black,
  citecolor=black,
  urlcolor=black,
  pdftitle={Distinction Reorganization and the Geometry of Understanding},
  pdfauthor={Flyxion}
}

\usepackage{amsmath, amssymb, amsthm}
\usepackage{mathtools}
\usepackage{enumitem}
\usepackage{booktabs}
\usepackage{array}
\usepackage{xcolor}
\usepackage{fancyhdr}
\setlength{\headheight}{15pt}

\pagestyle{fancy}
\fancyhf{}
\fancyhead[L]{\small\textit{Flyxion}}
\fancyhead[R]{\small\textit{Distinction Reorganization and the Geometry of Understanding}}
\fancyfoot[C]{\thepage}
\renewcommand{\headrulewidth}{0.3pt}

% Theorem environments
\theoremstyle{plain}
\newtheorem{theorem}{Theorem}[section]
\newtheorem{proposition}[theorem]{Proposition}
\newtheorem{lemma}[theorem]{Lemma}
\newtheorem{corollary}[theorem]{Corollary}
\theoremstyle{definition}
\newtheorem{definition}[theorem]{Definition}
\newtheorem{example}[theorem]{Example}
\theoremstyle{remark}
\newtheorem{remark}[theorem]{Remark}
\newtheorem{principle}[theorem]{Principle}

% Notation
\newcommand{\Dist}{\mathcal{D}}
\newcommand{\Adm}{\mathcal{A}}
\newcommand{\Reach}{\mathcal{R}}
\newcommand{\Proj}{\pi}
\newcommand{\Obs}{\mathcal{O}}
\newcommand{\PhiF}{\varPhi}
\newcommand{\PhiFop}{\varPhi_{\mathrm{op}}}
\newcommand{\PhiFid}{\varPhi^{*}}
\newcommand{\Tlib}{\mathcal{T}}
\newcommand{\Mod}{\mathcal{M}}
\newcommand{\Lam}{\Lambda}
\newcommand{\dA}{D_{\!\mathcal{A}}}
\newcommand{\nstar}{n^{*}}
\newcommand{\DG}{\mathcal{G}_{\Dist}}
\DeclareMathOperator{\Hent}{H}
\DeclareMathOperator{\Info}{I}

\title{%
  \LARGE\bfseries Distinction Reorganization and the\\
  Geometry of Understanding%
}

\author{Flyxion}

\date{2026}

\begin{document}

\maketitle
\thispagestyle{fancy}

\begin{abstract}
\noindent
Scientific revolutions, memory reconsolidation, archive rematching,
representation learning, mechanistic interpretability, and ontology repair
are customarily treated as distinct problems belonging to different
disciplines. This paper argues that they are manifestations of a single
underlying process: the reorganization of admissible distinctions into
a structure with greater reachability and recoverability.
The central claim is that understanding is not the accumulation of
descriptions but a \emph{trajectory} of distinction reorganization:
a sequence of structures $\Dist_0 \to \Dist_1 \to \cdots$ in which
each transition expands reachability, reduces residual density, and
retroactively remaps the accumulated observational archive.
We introduce a formal framework --- distinction geometry --- in which
distinctions are organized as dependency graphs, reachability is
defined via recovery operators, and residual structure density $\varPhi$
is given both an ideal formulation (via Kolmogorov complexity) and a
computable operational proxy (via distance to the learned admissibility
manifold).
The framework is used to formalize Kuhnian paradigm change as a
quantitative dynamical process ($\text{revolution} = \Delta\bar\varPhi \ll 0$),
to state the World-Model Crossover and Directionality Principle as
the formal boundary between observation-centered and model-centered
epistemology, to introduce the Representational Revision Principle
as a guard against conflating epistemological and ontological revision,
and to prove the Observational--Interventional Separation Theorem as
the formal foundation for the interpretability reframing.
The framework unifies six independent research programs ---
Frozen Processes, Persistent Anomalies, Waves of Collapse,
Admissibility Geometry, Distinguishability Geometry, and RSVP ---
as cross-sections of a single geometric object.
\end{abstract}

\newpage

\tableofcontents

\bigskip\hrule\bigskip

%% ============================================================
\section{Introduction}
\label{sec:intro}
%% ============================================================

There is a question that sits uncomfortably beneath almost every
contemporary discussion of artificial intelligence, cognitive science,
and the philosophy of science, yet is rarely stated directly:

\begin{quote}
\textit{What is the difference between a system that predicts correctly
and a system that understands?}
\end{quote}

The standard answer --- that understanding is reliable, general, or
robust prediction --- is demonstrably insufficient. A system may
achieve high predictive accuracy on every benchmark while collapsing
exactly the distinctions that would allow it to intervene, explain,
or generalize beyond its training distribution. This is not a failure
of capability. It is a failure of faithfulness: the representation
has discarded structure that matters.

This paper proposes a different answer. Understanding arises when a
representation preserves and reorganizes the \emph{admissible
distinctions} of a domain --- the distinctions whose collapse would
impair navigation, intervention, or explanation relative to the
purposes for which the representation exists. Understanding advances
not gradually but through \emph{collapse events}: discontinuous
reorganizations of the distinction structure that expand reachability,
enable rematching of prior observations, and increase the domain's
navigability.

This claim unifies a cluster of apparently unrelated phenomena.
Scientific revolutions are collapse events at civilizational scale.
Memory reconsolidation is collapse at the scale of individual belief.
Archive rematching is collapse in the domain of stored records.
Ontology repair in AI systems is collapse in the domain of learned
representations. Mechanistic interpretability is, properly understood,
the project of verifying whether collapse has occurred and what
structure it has revealed.

The framework rests on three core concepts.

\emph{Admissible distinctions} are those whose collapse would impair
performance on a given task or family of tasks. Admissibility is
always relative to a purpose; the same difference in the world may
be admissible for one task and inadmissible for another.

\emph{Reachability} is the capacity to express, recover, and operate
on a distinction within a given representation. A distinction that
exists in the world but is not reachable from a representation
does not belong to that representation's understanding.

\emph{Ontological deficit} is the frozen residue of collapsed
admissible distinctions: the structured anomalies that a
representation cannot compress, the persistent noise that resists
current vocabulary, the pressure that accumulates before a collapse
event releases it.

Together these concepts define a geometry --- the geometry of
understanding --- in which progress is not a smooth accumulation
but a sequence of reorganizations, each of which expands the space
of what is reachable.

The paper is structured as follows.
Section~\ref{sec:projection} establishes that every nontrivial
representation destroys distinctions (the Projection Theorem) and
introduces the core vocabulary.
Section~\ref{sec:admissibility} defines admissibility and ontological
deficit and distinguishes them from predictive failure.
Section~\ref{sec:phi} introduces residual structure density $\PhiF$
as the quantitative measure of accumulated distinction failure.
Section~\ref{sec:kuhn} reinterprets Kuhnian paradigm change as
distinction reorganization, making the sociology of science
geometrically precise.
Section~\ref{sec:crossover} states the World-Model Crossover Principle
and the Directionality Principle as the formal boundary between
observation-centered and model-centered epistemology.
Section~\ref{sec:memory} shows that human memory behaves exactly
as the framework predicts.
Section~\ref{sec:ai} applies the framework to AI interpretability.
Section~\ref{sec:unification} shows how the framework unifies
several independent research programs.
Section~\ref{sec:conclusion} states the master claim and its
corollaries.

%% ============================================================
\section{Projection and Distinction Destruction}
\label{sec:projection}
%% ============================================================

\subsection{Representations as projections}

Every representation is a map from a richer domain to a poorer one.
This is not a defect but a necessity: a map that preserved every
detail of the territory would be useless. The question is not whether
information is lost but which information.

\begin{definition}[Representation and projection]
\label{def:rep}
Let $\Omega$ be a domain of cases. A \emph{representation} is a
surjective map $\Proj : \Omega \to M$ where $M$ is the representational
space. We call $\Proj$ a \emph{projection} and $M$ the \emph{model space}.
\end{definition}

\begin{definition}[Distinction]
\label{def:dist}
A \emph{distinction} on $\Omega$ is a binary relation
$d \subseteq \Omega \times \Omega$ such that $d(x, y)$ holds when
$x$ and $y$ are treated differently for some purpose.
A distinction induces a partition of $\Omega$ into equivalence classes
under the complementary relation $\lnot d$.
\end{definition}

A flat set of such relations is too weak a substrate for the framework.
The concepts of strata, ceilings, and reachability all presuppose that
distinctions stand in dependency relations to one another: collapsing
a high-level distinction may destroy downstream ones without affecting
peers. We therefore organize distinctions into a richer structure.

\begin{definition}[Distinction dependency graph]
\label{def:ddg}
A \emph{distinction dependency graph} $\DG = (V, E)$ over domain $\Omega$
is a directed acyclic graph in which:
\begin{enumerate}[label=(\roman*)]
  \item Each vertex $v \in V$ corresponds to a distinction $d_v$ on $\Omega$.
  \item A directed edge $(u \to v) \in E$ indicates that $d_v$
        \emph{depends on} $d_u$: drawing $d_v$ presupposes that $d_u$
        has already been drawn. Collapsing $d_u$ therefore propagates
        upward, destroying all $d_v$ reachable from $u$ via directed paths.
\end{enumerate}
The \emph{distinction structure} $\Dist$ is the pair $(\Omega, \DG)$.
\end{definition}

\begin{example}[Biological taxonomy as a distinction graph]
Consider the distinctions: \textit{living thing} ($d_1$),
\textit{animal} ($d_2$), \textit{vertebrate} ($d_3$),
\textit{mammal} ($d_4$), \textit{dog} ($d_5$).
The dependency graph has edges $d_1 \to d_2 \to d_3 \to d_4 \to d_5$.
Collapsing $d_4$ (the mammal distinction) propagates upward and destroys
$d_5$, but leaves $d_1, d_2, d_3$ intact.
This asymmetry --- which the flat binary-relation definition cannot
represent --- is what gives distinction structures their geometry.
\end{example}

\begin{definition}[Graph reachability and collapse]
\label{def:collapse}
A distinction $d_v$ is \emph{collapsed} by $\Proj$ if there exist
$x, y \in \Omega$ with $d_v(x,y)$ but $\Proj(x) = \Proj(y)$.

A distinction $d_v$ is \emph{graph-reachable} in $\DG$ from a set
$S \subseteq V$ if there exists a directed path from some $u \in S$
to $v$ in $\DG$. Collapsing a distinction $d_u$ renders
graph-unreachable all distinctions depending on $d_u$.
\end{definition}

The following result is elementary but foundational.

\begin{theorem}[Projection Theorem]
\label{thm:projection}
Let $\Proj : \Omega \to M$ be any nontrivial projection
(i.e., $|M| < |\Omega|$). Then $\Proj$ collapses at least one
distinction on $\Omega$.
\end{theorem}

\begin{proof}
Since $|M| < |\Omega|$, by the pigeonhole principle there exist
distinct $x, y \in \Omega$ with $\Proj(x) = \Proj(y)$.
The pair $(x,y)$ constitutes a collapsed distinction.
\end{proof}

The Projection Theorem is not a theorem about compression
in the technical sense; it is a theorem about representation in general.
Every simplification, abstraction, categorization, and model
necessarily collapses some distinctions. The question of epistemology
is which distinctions are collapsed and which are preserved.

\subsection{The Weierstrass obstruction}

A tacit assumption in many representational frameworks is that
a sufficiently rich finite vocabulary can approximate any signal
to arbitrary precision. This assumption fails for signals with
persistent fine-scale structure.

\begin{remark}[Weierstrass obstruction]
\label{rem:weierstrass}
The Weierstrass function $f(x) = \sum_{n=0}^\infty a^n \cos(b^n \pi x)$
(with $0 < a < 1$, $b$ odd, $ab > 1 + \frac{3\pi}{2}$) is continuous
everywhere and differentiable nowhere. No finite library of smooth
templates achieves finite residual on such a signal: new oscillations
appear at every scale of magnification.

For a representation whose template library contains only smooth
primitives, the Weierstrass-type signal constitutes a permanent
ontological deficit. The deficit is not resolvable by adding more
smooth templates; it requires the admission of a qualitatively
different template kind (recursive, scale-coupling templates).
\end{remark}

This observation is more than a mathematical curiosity.
It establishes that ontological deficit can be \emph{structural}
rather than merely quantitative: not a matter of having too few
templates, but of having the wrong kind.

%% ============================================================
\section{Admissibility and Ontological Deficit}
\label{sec:admissibility}
%% ============================================================

\subsection{Admissible distinctions}

The Projection Theorem tells us that distinction collapse is
inevitable. The theory of admissibility identifies which collapses
constitute failures.

\begin{definition}[Admissibility]
\label{def:adm}
Let $\tau$ be a task or purpose. A distinction $d_v \in V$ is
\emph{admissible for $\tau$} if collapsing $d_v$ would impair
performance on $\tau$.
The \emph{admissibility manifold} $\Adm_\tau \subseteq V$ is the set
of all distinctions admissible for $\tau$.
\end{definition}

Before developing the consequences of admissibility, we supply the
formal definition of reachability that the framework requires throughout.

\begin{definition}[Reachability]
\label{def:reach}
A distinction $d_v$ is \emph{reachable} from a representation
$\Proj : \Omega \to M$ if there exists a computable function
$R_v : M \to \{0,1\}$ such that $R_v(\Proj(x)) = \mathbf{1}[d_v(x,\cdot)]$
for all $x$ in an appropriate reference class.
We call $R_v$ a \emph{recovery operator} for $d_v$.

The \emph{reachability set} of $\Proj$ relative to task $\tau$ is
\[
  \Reach(\Proj, \tau)
  \;=\;
  \bigl\{d_v \in \Adm_\tau : \exists\, R_v \text{ as above}\bigr\}.
\]
A representation \emph{understands} domain $\Omega$ for task $\tau$
to the degree that $\Adm_\tau \subseteq \Reach(\Proj, \tau)$.
\end{definition}

Reachability, so defined, is the formal counterpart of the Witness
Principle introduced in Section~\ref{sec:ai}: a distinction is
reachable precisely when a recovery operator witnesses it.
The reachability set is therefore the set of distinctions for which
constructive evidence exists.

Several features of admissibility deserve immediate emphasis.

First, admissibility is \emph{task-relative}. The distinction between
two chemical compounds that are pharmacologically identical but
isotopically distinct is admissible for nuclear physics but
inadmissible for most pharmacological tasks. This relativity is
not a defect of the framework but a constitutive feature of
understanding: understanding is always understanding \emph{for some
purpose}.

Second, the admissibility manifold has a natural stratification.
Some distinctions are admissible for all tasks in a domain
(structural distinctions); others are admissible only for specific
interventions (functional distinctions); others only for counterfactual
reasoning (causal distinctions). These strata correspond to
different levels of understanding.

Third, admissibility is not observable from outside the task.
A representation may appear faithful on observational tasks while
collapsing distinctions required for interventional or causal tasks.
This is the source of the prediction trap.

\begin{definition}[Projection failure]
\label{def:projfail}
A projection $\Proj$ \emph{fails} relative to task $\tau$ if it
collapses any distinction in $\Adm_\tau$. The \emph{admissibility
distortion} is
\[
  \dA(\Proj, \tau)
  \;=\;
  \Pr_{(x,y) \in \Adm_\tau}\!\bigl[\Proj(x) = \Proj(y)\bigr].
\]
A representation with $\dA = 0$ is \emph{faithful} for $\tau$;
with $\dA > 0$ it has a \emph{projection failure}.
\end{definition}

\subsection{Ontological deficit}

Ontological deficit is a stronger condition than projection failure.
It arises when the representational vocabulary is incapable of
expressing a distinction at all, not merely when a distinction
has been collapsed for a particular input.

\begin{definition}[Ontological deficit]
\label{def:ontdef}
A representation $\Proj$ has an \emph{ontological deficit} with
respect to distinction $d$ if no state of $M$ corresponds to
drawing the distinction $d$: the distinction is not merely
collapsed but inexpressible.

Formally, there is no partition of $M$ that represents $d$,
even approximately.
\end{definition}

\begin{example}[Pre-Darwinian natural history]
Pre-Darwinian taxonomy drew distinctions between species by
morphological similarity. The distinction between homologous and
analogous structures --- structures sharing a common ancestor
versus structures convergently evolved --- was not merely collapsed
by pre-Darwinian taxonomy but inexpressible within it, because
the vocabulary of common descent did not exist.
This is an ontological deficit, not a projection failure.
\end{example}

\begin{example}[Language model token prediction]
A language model trained solely on next-token prediction may
achieve low perplexity while having an ontological deficit with
respect to the distinction between \emph{true} and \emph{false
but fluent} continuations. The distinction is not represented
in the model's output space if no training signal distinguished
it; no linear probe on the model's internal representations
will reliably recover it.
\end{example}

The practical significance of ontological deficit is that it is
\emph{invisible from within}. A framework with an ontological
deficit cannot, using its own resources, detect the deficit.
It experiences the consequences --- persistent anomalies, residual
structure that resists compression, asymmetric prediction errors ---
but cannot name their source. The deficit becomes visible only
from the perspective of a successor framework that has admitted
the missing distinction.

%% ============================================================
\section{Residual Structure Density and the Pressure Toward Collapse}
\label{sec:phi}
%% ============================================================

\subsection{Definition and properties}

Ontological deficit is a qualitative description of representational
failure. The following construction provides a quantitative measure.

\begin{definition}[Residual structure density: ideal and operational]
\label{def:phi}
Let $\Proj : \Omega \to M$ be a representation and $\Tlib$ the current
template library.

The \emph{ideal residual density} at $x \in \Omega$ is
\[
  \PhiFid(x) \;=\; K(x \mid \Tlib) \;-\; K(x \mid \Tlib^*)
\]
where $K(\cdot \mid \cdot)$ is conditional Kolmogorov complexity and
$\Tlib^*$ is the library achieving minimal description length.
$\PhiFid$ measures the theoretical gap between current and optimal
understanding; it is the benchmark against which progress is assessed.

The \emph{operational residual density} at $x \in \Omega$ is
\[
  \PhiFop(x) \;=\; \min_{t \in \Tlib}\, d_{\Adm}(x,\, \mathcal{M}_t)
\]
where $\mathcal{M}_t$ is the manifold generated by template $t$ and
$d_{\Adm}$ is a metric on the admissibility space.
$\PhiFop$ is the distance from $x$ to the nearest template boundary:
it is computable, measurable, and directly actionable.

We write $\PhiF$ when referring to either form.
The total residual is $\bar\PhiF = \int_\Omega \PhiF(x)\,d\mu(x)$.
\end{definition}

The two quantities serve complementary roles. $\PhiFid$ provides the
conceptual foundation: it is what the framework is ultimately measuring.
$\PhiFop$ provides the research instrument: it is what practitioners
can compute as a distance from data points to the learned manifold,
and its gradient $\nabla\PhiFop$ gives the direction in which new
template admissions are most likely to be productive.
This separation --- between the theoretical ideal and the operational
proxy --- is standard practice in information-theoretic frameworks
and insulates the $\PhiF$ concept from the uncomputability objection
while retaining its explanatory force.

Three properties hold for both forms.

First, $\PhiF \geq 0$ everywhere, with equality if and only if
the current vocabulary is optimal for $x$.

Second, $\PhiF$ is \emph{structured}: it is concentrated at points
where the representation fails. Structured concentration of $\PhiF$
(systematic residual rather than random noise) is the operational
signature of ontological deficit. Random residuals do not define a
gradient; structured ones do.

Third, $\PhiF$ exerts \emph{pressure} on the representation.
The gradient $\nabla\PhiFop$ defines the direction in which new template
admissions are likely to reduce $\bar\PhiF$. The anomaly field is the
map of the investigation frontier.

\subsection{The pressure toward collapse}

A template library $\Tlib$ under pressure from a concentrated
$\PhiF$ field is in an unstable equilibrium. New templates can be
admitted --- speculatively, on credit --- if they reduce the total
description length:
\[
  \Lam(\Tlib) \;=\; L(\Tlib) \;+\; \sum_{x \in \Obs} L(x \mid \Tlib)
\]
where $L(\Tlib)$ is the cost of the library itself (model complexity)
and $\sum L(x \mid \Tlib)$ is the total residual cost.

A template $t$ is \emph{admitted} if $\Lam(\Tlib \cup \{t\}) < \Lam(\Tlib)$.
This Minimum Description Length criterion is the formal analog
of Ockham's razor: templates earn their place by explaining more
than they cost.

When a successful template is admitted, the $\PhiF$ field undergoes
a rapid transition: regions of high residual that the new template
addresses are resolved into the crystallized entropy field $S$.
The drop in $\bar\PhiF$ is discontinuous if the template addresses
many observations simultaneously.

This discontinuity is the collapse event.

%% ============================================================
\section{Paradigm Change as Distinction Reorganization}
\label{sec:kuhn}
%% ============================================================

\subsection{The Kuhnian framework revisited}

Kuhn's account of scientific revolutions~\cite{kuhn1962} identifies
a pattern: long periods of normal science, punctuated by crises
and revolutionary reorganizations. This account is persuasive as
history but resists formalization. What exactly is a paradigm?
Why do anomalies accumulate? What makes a revolution rather than
a refinement?

The distinction-geometry framework provides precise answers.
The translation below is not merely metaphorical; each entry on the
right admits mathematical treatment.

\begin{definition}[Kuhnian translation]
\label{def:kuhn}
Under the distinction-geometry framework:
\begin{align*}
  \text{paradigm} &\;\longleftrightarrow\; \text{distinction structure } \Dist = (\Omega, \DG) \\
  \text{normal science} &\;\longleftrightarrow\; \text{local optimization within } \Dist \\
  \text{anomaly} &\;\longleftrightarrow\; \text{residual structure density } \PhiF \\
  \text{anomaly pressure} &\;\longleftrightarrow\; \bar\PhiF \\
  \text{crisis} &\;\longleftrightarrow\; \text{structured, intractable } \PhiF \\
  \text{revolution} &\;\longleftrightarrow\; \text{distinction reorganization (collapse event)} \\
  \text{Kuhn loss} &\;\longleftrightarrow\; \text{distinction retirement} \\
  \text{incommensurability} &\;\longleftrightarrow\; \text{non-overlapping admissibility manifolds}
\end{align*}
\end{definition}

The translation makes Kuhnian dynamics quantitative in the following
precise sense.

\begin{proposition}[Quantitative Kuhn]
\label{prop:kuhn-quant}
Under the distinction-geometry framework:
\begin{enumerate}[label=(\roman*)]
  \item \emph{Anomaly pressure} is measured by total residual density:
        $\text{pressure} = \bar\PhiF$.
  \item \emph{Normal science} is operation in a regime where
        $\frac{d\bar\PhiF}{dt} \approx 0$ under routine template additions.
  \item \emph{Crisis} occurs when $\bar\PhiF$ is both large and structured:
        residuals form a nonzero gradient $\nabla\PhiFop \neq 0$.
  \item \emph{Revolution} is a collapse event characterized by
        $\Delta\bar\PhiF \ll 0$: a discontinuous drop in total residual
        following the admission of a new distinction $d^*$.
\end{enumerate}
\end{proposition}

Whether or not the formalism is fully computable in historical cases,
the conceptual transformation is important: paradigm change ceases to
be a historical narrative and becomes a dynamical process with
measurable precursors (structured $\PhiF$) and measurable outcomes
($\Delta\bar\PhiF$). The sociology of science becomes, in principle,
physics of information.

\subsection{Normal science as local optimization}

Within a fixed distinction structure $\Dist$, the archive performs
routine encoding: new observations are matched to existing templates,
residuals are minimized within the current vocabulary, and the
reachability set $\Reach(\rho, \tau)$ expands incrementally.
This is the operational definition of normal science: progress
that does not require new distinctions, only better use of
existing ones.

Normal science is not stagnant. Within a stratum, significant
progress is possible. But every stratum has a ceiling: the
structural barrier defined by the distinctions it cannot express.
The ceiling becomes visible when $\PhiF$ accumulates in ways that
resist normal-science compression --- when anomalies are not
random but systematic.

\begin{definition}[Stratum ceiling]
The \emph{ceiling} of a distinction structure $\Dist$ is the
minimum description length achievable within $\Dist$:
\[
  \Lam^*(\Dist) = \min_{\Tlib \subseteq \Dist} \Lam(\Tlib).
\]
The ceiling is positive if $\Dist$ has an ontological deficit
relative to the data-generating process.
\end{definition}

\subsection{Anomaly accumulation as pressure}

Anomalies in Kuhn's account are observations that the paradigm
cannot explain. In the distinction-geometry framework, they are
regions of high, structured $\PhiF$: observations that resist
compression under the current template library.

The crucial word is \emph{structured}. Random residuals are
uninformative. Structured residuals --- residuals that have a
pattern, that vary systematically with known variables, that
cluster in the same regions across independent observers ---
carry information about the missing distinction. They define
the misfit gradient: the direction in which new template
admissions are likely to be productive.

Anomaly accumulation is therefore not a pathology but an
information-bearing process. The archive's persistent anomalies
are its research program: they are the pressure that drives
collapse.

\subsection{Revolution as distinction reorganization}

A scientific revolution is a collapse event at civilizational
scale: a transition $\Dist \to \Dist'$ in which a new distinction
$d^*$ reorganizes the entire distinction structure, reduces
$\bar\PhiF$ discontinuously, and enables rematching of the
accumulated observational archive.

\begin{proposition}[Structure of a revolution]
A distinction reorganization event $\Dist \to \Dist'$ driven by
the admission of $d^*$ has the following structure:
\begin{enumerate}[label=(\roman*)]
  \item \textbf{Inexpressibility phase}: $d^*$ is not reachable
        in $\Dist$; its consequences appear as structured $\PhiF$.
  \item \textbf{Crystallization phase}: $d^*$ becomes expressible
        through new vocabulary, instrumentation, or formal notation.
        At this stage $d^*$ is not yet integrated into $\Dist$.
  \item \textbf{Reorganization phase}: The admission of $d^*$
        retroactively restructures $\Dist$. Cases previously
        classified together are separated; cases previously
        separate are unified.
        $\bar\PhiF$ drops discontinuously.
  \item \textbf{Rematching phase}: The observational archive
        is reinterpreted under $\Dist'$. Old observations yield
        new descriptions, often with greater precision than
        was possible at the time of observation.
\end{enumerate}
\end{proposition}

\begin{example}[Caloric to thermodynamic entropy]
\label{ex:caloric}
The caloric theory treated \emph{amount of caloric} as an admissible
distinction. Thermodynamics retired this distinction and admitted
\emph{entropy} (the distinction between states with many and few
microstates). The reorganization collapsed the caloric distinction
(reachable in the old stratum but retired in the new) while
admitting the entropy distinction (inexpressible in the old stratum).

The rematching phase was extensive: steam engine efficiency data,
previously described using caloric quantities, was reinterpreted
in entropy terms, yielding Carnot's efficiency theorem and its
consequences for what is and is not possible for heat engines.
The observational archive had not changed; the distinction structure
had reorganized.
\end{example}

\subsection{Kuhn loss and distinction retirement}

Kuhn loss is the observation that something may be expressible
within an old paradigm that is not expressible within its
successor. In the distinction-geometry framework, Kuhn loss is
distinction retirement: a distinction $d$ that was part of $\Dist$
may not appear in $\Dist'$.

This is not symmetric with distinction admission. A retired
distinction was real within the old framework --- it supported
genuine inferences and successful manipulations --- but the new
framework explains those successes without it, or explains them
as approximations of a more fundamental distinction.

Kuhn loss does not refute the progress account of science.
It refines it: progress is the growth of the reachability set
over the \emph{admissible} distinctions, which may require retiring
distinctions that were previously admitted but turn out to be
approximations, confabulations, or ontological artifacts of the
old vocabulary.

\subsection{Incommensurability without relativism}

Kuhn's incommensurability thesis has often been read as implying
that paradigms cannot be compared and that scientific revolutions
have no rational basis. The distinction-geometry framework
dissolves this difficulty.

\begin{principle}[Local incommensurability, global comparability]
\label{prin:incomm}
Two distinction structures $\Dist$ and $\Dist'$ are
\emph{locally incommensurable} on a distinction $d$ if
$d \in \Dist$ but drawing $d$ is inconsistent with the
organizing principles of $\Dist'$ (or vice versa).

But $\Dist$ and $\Dist'$ are \emph{globally comparable}
if their admissibility manifolds overlap on a sufficiently
large common ground.

Rational evaluation across a revolution proceeds via the common
ground: observations expressible in both frameworks, tasks
performable under both distinction structures, predictions
where the structures agree or disagree in testable ways.
\end{principle}

Incommensurability is local; comparability is global.
This is why revolutions are difficult (locally, the new vocabulary
is alien) but not irrational (globally, the new structure explains
more of the shared observations and reduces $\bar\PhiF$).

%% ============================================================
\section{The World-Model Crossover Theorem}
\label{sec:crossover}
%% ============================================================

\subsection{The crossover condition}

Every representation starts as a summary of observations.
Initially, observations organize the model: each new datum updates
the template library, and the model's content is determined by
what it has seen. There is, however, a threshold beyond which
the model's internal structure becomes more informative about
the domain than any individual observation.

This is the World-Model Crossover.

\begin{principle}[World-Model Crossover]
\label{prin:crossover}
Let $\Mod$ be a representation of domain $\Omega$ built from a
sequence of observations $\{O_1, O_2, \ldots, O_n\}$.
Define the \emph{crossover threshold} $\nstar$ as the smallest
$n$ such that
\[
  \Info(\Mod_n;\, \Omega) \;>\; \max_{i \leq n}\, \Info(O_i;\, \Omega)
\]
where $\Info(\cdot\,;\,\cdot)$ denotes mutual information.

For $n < \nstar$, the observations are the primary information-bearing
objects about $\Omega$; the model is their summary.

For $n \geq \nstar$, the model is the primary information-bearing
object; individual observations are perturbations to the model.

An equivalent formulation in the distinction-geometry framework:
$\nstar$ is the point at which $\Reach(\Proj, \tau)$ becomes
independent of the size and recency of the training dataset ---
the model's internal geometry, rather than any individual observation,
provides the best guide for future inference.
\end{principle}

\begin{remark}
The mutual-information inequality is an informative criterion but
not yet a theorem with a constructive proof in the general setting.
The formal crossover condition remains a conjecture in full generality.
Its value is explanatory: it identifies the qualitative transition
and defines what would need to be established for the claim to be proved.
\end{remark}

The crossover is not merely quantitative. It marks a \emph{qualitative}
transition in the epistemological role of the model.

\begin{principle}[Directionality Principle]
\label{prin:direction}
Below $\nstar$:
\[
  \text{observations} \;\longrightarrow\; \text{model}
  \quad\text{(observations organize theories)}
\]
Above $\nstar$:
\[
  \text{model} \;\longrightarrow\; \text{observations}
  \quad\text{(theories organize observations)}
\]
The directionality is informational, not causal.
Below the crossover, an observation is primarily \emph{data}:
something to be explained by the model.
Above the crossover, an observation is primarily \emph{evidence}:
something that confirms, disconfirms, or anomalously pressures
the model, adding to the $\PhiF$ field if it cannot be absorbed.
\end{principle}

The Directionality Principle is the epistemologically significant
claim. The inequality in Principle~\ref{prin:crossover} is its
formal vehicle; the principle is what readers will recognize in
the actual practice of mature sciences. Newtonian mechanics, once
established, did not wait for further pendulum swings to organize
its claims about motion. It used its own internal structure to predict
pendulum behavior, and individual pendulum swings became tests of
the theory rather than its constituents. The crossover had occurred.

This transition defines the boundary between data collection and
theory formation, between empirical science and theoretical science,
between an archive and a world-model.

\subsection{The byte as a unit of surprise}

The crossover has a striking operational consequence for storage.

Below $\nstar$, the cost of storing new observations tracks the
\emph{data rate} of the sensors: more observations, more storage.

Above $\nstar$, the cost tracks the \emph{novelty rate} of the world:
\[
  \text{storage cost} \;\approx\; \Hent(O \mid \Mod)
\]
where $\Hent(O \mid \Mod)$ is the entropy of an observation
\emph{given the model}. If the model already knows what the
sensors will record, the storage cost approaches zero.
A 16K camera recording an empty, unchanging room costs almost
nothing to store above $\nstar$, because the model already
knows the room.

The byte is no longer a unit of measurement; it is a unit of
\emph{surprise}. The archive spends bits only on what its
model could not predict.

\subsection{Crossover and the diversity requirement}

The crossover threshold $\nstar$ is not reached by accumulating
more observations of the same type. It requires \emph{diversity
of kernels}: observations from different measurement modalities
that triangulate the same underlying domain from different angles.

\begin{proposition}[Diversity requirement]
\label{prop:diversity}
Let $\{K_1, K_2, \ldots, K_m\}$ be measurement kernels
(e.g., visual, acoustic, geometric).
The joint information $\Info(\{K_i\};\, \Omega)$ exceeds
$\max_i \Info(K_i;\, \Omega)$ when the kernels expose
orthogonal aspects of the domain.

Formally, $\nstar$ is minimized by maximizing the diversity
of the kernel family:
\[
  \nstar \propto \frac{1}{\sum_{i \neq j} \Info(K_i; K_j \mid \Omega)}
\]
Highly correlated kernels (many cameras from similar angles)
delay the crossover; orthogonal kernels (visual plus acoustic
plus structural) accelerate it.
\end{proposition}

No single modality reaches the crossover alone. Genuine world-models
emerge from the integration of multiple measurement kinds, because
no single kind exposes the full structure of the domain.

This is not merely a technical observation. It explains why
scientific understanding depends on the coordination of multiple
experimental traditions, and why disciplines that share data
across modalities progress faster than those that do not.

\subsection{Post-crossover rematching}

Above the crossover, a new observation that conflicts with the
model does not merely update the model's beliefs. It generates
an anomaly in the $\PhiF$ field, creating pressure for a
reorganization.

The post-crossover archive therefore exhibits the same dynamics
as the pre-revolution scientific community in Kuhn's account:
anomalies accumulate, misfit gradients point toward missing
distinctions, and collapse events are triggered by template
admissions that resolve the pressure.

The crossover is not a permanent state but a moving threshold.
Each collapse event raises the model's informational content,
but also opens new strata of the domain that the model cannot
yet address. The $\PhiF$ field never reaches zero; it reorganizes.
Understanding is not a destination but a process of continual
reorganization.

%% ============================================================
\section{Memory as Distinction Reorganization}
\label{sec:memory}
%% ============================================================

\subsection{The recording illusion}

The folk model of memory presents recollection as retrieval: a
record stored at encoding is located and returned intact.
This picture is not merely imprecise but structurally wrong in
a way that the distinction-geometry framework clarifies.

A representation stores not observations but \emph{the distinctions
an observation supported}. Recollection is reconstruction from
those distinctions as they are organized in the current
distinction structure. Because the current structure differs from
the structure at encoding, recollection is necessarily reconstruction,
not retrieval.

\begin{principle}[Memory as reconstruction]
\label{prin:memory}
Recollection of event $e$ is the inference to the best
distinction-preserving completion of the current distinction
structure $\Dist_t$ consistent with the partial record
encoded at time of observation.

Since $\Dist_t \neq \Dist_{t_0}$ in general, the recollection
is filtered through distinctions that did not exist at encoding.
This is not a defect but a constitutive feature: the past is
rendered in the vocabulary of the present.
\end{principle}

The phrase ``rendered in the vocabulary of the present'' requires a
guard against a common misreading.

\begin{principle}[Representational Revision Principle]
\label{prin:reprrev}
Distinction reorganization modifies the \emph{representation}
of prior observations, not the observations themselves.

Formally: a collapse event $\Dist \to \Dist'$ changes the encoding
$\Proj_{t'}(e) \neq \Proj_t(e)$ for observations $e$ encoded
before the event, but the observation $e \in \Omega$ itself is
unchanged.

The past is not rewritten. The representation of the past is rewritten.
\end{principle}

This principle insulates the framework from the objection that it
conflates epistemology with ontology --- an objection that has been
leveled against strong readings of the free-energy principle
and related Markov blanket formalisms. The distinction-geometry
framework is entirely epistemological: it concerns how observations
are encoded and recovered, not what the observations themselves are.
The revision of a memory under reconsolidation changes the \emph{render},
not the event.

This prediction is well confirmed. Bartlett's experiments
on reconstructive memory~\cite{bartlett1932}, the misinformation
effect studied by Loftus~\cite{loftus1979}, and the
reconsolidation research of Nader and colleagues~\cite{nader2000}
all document exactly the behavior the framework predicts:
recall is shaped by post-encoding knowledge, schema, and belief.

\subsection{Reconsolidation as local collapse}

Memory reconsolidation --- the phenomenon in which retrieved
memories become temporarily labile and subject to modification ---
is, in the distinction-geometry framework, a local collapse event.

When a memory is retrieved, it is reconstructed within the current
distinction structure $\Dist_t$. If $\Dist_t$ contains distinctions
that were absent at encoding time $t_0$, the reconstruction may
differ from the original. The reconsolidation window is the period
during which the reconstructed memory is re-encoded using the
current (potentially reorganized) vocabulary.

This is not distortion in any pejorative sense. It is the
faithful application of the distinction-reorganization process
to the domain of episodic memory. The memory becomes more
\emph{reachable} --- more efficiently represented within the
current distinction structure --- at the cost of some fidelity
to the original encoding.

Reconsolidation is therefore not a vulnerability of the memory
system but evidence of its design: a rematching archive does not
store frozen records but living theories, and living theories
must be updated when the vocabulary improves.

\subsection{The compressibility of the past}

A counterintuitive prediction of the framework is that the past
\emph{becomes more compressible over time} as the distinction
structure matures.

Early in learning, autobiographical memories are encoded in
coarse-grained vocabularies: the memory of a trip to Paris at
age twelve is encoded in whatever distinction structure was
available at age twelve. As the distinction structure matures
--- as new templates are admitted, as more subtle distinctions
become reachable --- the same event can be re-encoded more
efficiently. Details that were stored as raw impressions become
pointers to general templates; patterns that were opaque become
recognizable instances of known categories.

This is what we experience as the wisdom of retrospection:
the sense that past experiences become more \emph{intelligible}
as we mature, not because the events changed but because our
distinction structure has reorganized to make them more compressible.

%% ============================================================
\section{Interpretability as Distinction Verification}
\label{sec:ai}
%% ============================================================

\subsection{The interpretability reframing}

Mechanistic interpretability in AI research is commonly framed
as the project of understanding what is happening inside a
trained model: identifying circuits, features, attention patterns,
and internal representations that explain the model's behavior.

The distinction-geometry framework reframes this project.
Interpretability is not primarily about \emph{describing}
internal structure. It is about \emph{verifying} whether the
model has preserved the admissible distinctions required for
navigation, explanation, and intervention.

A description of internal structure that does not answer this
question is a description, not an explanation.

\begin{principle}[Interpretability as distinction verification]
\label{prin:interp}
A mechanistic interpretability claim is \emph{explanatory}
only if it demonstrates that the identified internal structure
preserves a specific admissible distinction.

An interpretability claim that identifies a feature without
showing how that feature preserves a distinction in
$\Adm_\tau$ for some task $\tau$ is a description of the
representation, not an explanation of the understanding.
\end{principle}

This principle has sharp practical consequences.

\subsection{Chain-of-thought and the trajectory question}

Large language models prompted to produce chain-of-thought
reasoning exhibit markedly improved performance on tasks
requiring multi-step inference. The standard interpretation
is that chain-of-thought forces the model to externalize
its reasoning trajectory, making previously internal steps
visible and correctable.

The distinction-geometry framework raises a more pointed question:

\begin{quote}
Is the model \emph{performing} the reasoning trajectory,
or merely \emph{describing} a trajectory it did not use?
\end{quote}

These are not the same. A model could produce a fluent verbal
description of a reasoning process that is entirely post-hoc
confabulation, while the actual computation that produced the
answer preserved no admissible distinctions from the described
steps.

The test is not whether the chain-of-thought is coherent but
whether the distinctions drawn in the chain-of-thought
are preserved in the computation. If removing a step from the
chain changes the model's output in ways that track the causal
role of that step, the chain-of-thought preserves admissible
distinctions. If the chain is rearranged or corrupted without
affecting the output, it does not.

\subsection{The recoverability requirement}

A stronger condition than distinction preservation is
\emph{distinction recoverability}: the ability to extract
a specific distinction from the representation using a
designated recovery operator.

\begin{definition}[Recoverability]
\label{def:recover}
A distinction $d_v$ is \emph{recoverable} from representation $\Proj$
via operator $R_v$ if there exists a map $R_v : M \to \{0,1\}$ such that
$R_v(\Proj(x)) = \mathbf{1}[d_v(x,\cdot)]$ for all $x$ in the relevant domain.
\end{definition}

\begin{principle}[Witness principle]
\label{prin:witness}
A claim that distinction $d_v$ is represented in a model $M$
requires a recovery operator $R_v$ as constructive evidence.
Without a recovery operator, the claim is an assertion, not a proof.
\end{principle}

The witness principle is deliberately stringent. It excludes
post-hoc attribution stories that are consistent with the
model's behavior but do not reflect a genuine distinction
in the model's internal structure.

This is not a demand for perfect recovery. Approximate
recoverability --- with known error bounds --- is sufficient.
But the recovery operator must exist and be demonstrated,
not merely hypothesized.

\subsection{Observational and interventional separation}

The admissibility framework distinguishes two levels of distinction:
those separating cases that look different (\emph{observational})
and those separating cases that respond differently under intervention
(\emph{interventional}). A representation that is faithful at the
observational level may fail at the interventional level.

\begin{theorem}[Observational--Interventional Separation]
\label{thm:oisep}
Let $\Proj : \Omega \to M$ be a representation and let $\tau_{obs}$,
$\tau_{int}$ denote the observational and interventional task families.
Then:
\[
  \dA(\Proj, \tau_{obs}) = 0
  \;\not\Rightarrow\;
  \dA(\Proj, \tau_{int}) = 0.
\]
Observational faithfulness does not entail interventional faithfulness.
Equivalently, there exist cases $x, y \in \Omega$ such that
$\Proj(x) = \Proj(y)$ (observationally identical) but $do(x) \neq do(y)$
(interventionally distinct), where $do(\cdot)$ denotes the response
to a perturbation.
\end{theorem}

\begin{proof}
It suffices to exhibit a representation that collapses an interventional
distinction while preserving all observational ones. Let
$\Omega = \{x_1, x_2, x_3\}$ where $x_1, x_2$ have identical
observable properties but distinct responses to perturbation $p$:
$do_p(x_1) = r_1 \neq r_2 = do_p(x_2)$, and $x_3$ is observationally
distinct. Define $\Proj(x_1) = \Proj(x_2) = m_1$, $\Proj(x_3) = m_2$.
Then $\dA(\Proj, \tau_{obs}) = 0$ (only observationally distinct cases
are separated) but $\dA(\Proj, \tau_{int}) > 0$ (the interventional
distinction $d(x_1, x_2)$ is collapsed).
\end{proof}

The theorem is a formal statement of a familiar intuition in the
causal inference literature~\cite{pearl2009}: correlation does not
imply intervention-reliability. Its significance here is that it
gives the interpretability reframing a formal result to anchor on.
A model that achieves low test error may be intervening on a
statistical proxy rather than the genuine causal variable; the
Observational--Interventional Separation Theorem says this is
not merely possible but structurally expected when the training
signal is purely observational.

The practical test is: if the model's claimed internal distinction
$d_v$ is manipulated directly (via causal intervention on the
representation), does the output change in ways that track
the real-world causal role of $d_v$? If not, the distinction has
not been recovered; it has been described.

%% ============================================================
\section{Unification: Six Programs, One Geometry}
\label{sec:unification}
%% ============================================================

The framework developed in this paper was not designed to unify
independent research programs. The unification is a consequence
of the framework's generality. Nevertheless, it is worth making
the connections explicit.

\medskip
\noindent\textbf{Frozen Processes.}
The Frozen Processes program~\cite{flyxion2026b} studies what happens
when a distinction structure is assumed fixed: $\Adm_t = \Adm_0$.
In the language of the present paper, a frozen-process assumption
blocks distinction reorganization. The consequences are exactly what
Frozen Processes documents: the accumulation of $\PhiF$ in the form
of unexplained anomalies, the inability to rematch historical
observations, and the eventual crisis that forces a reorganization.
Frozen Processes is the study of the \emph{obstruction to collapse events}.

\medskip
\noindent\textbf{Persistent Anomalies.}
The Persistent Anomalies program studies structured residuals that
resist compression under the current template library.
In the present framework, a persistent anomaly is a concentrated,
structured region of the $\PhiF$ field whose misfit gradient points
toward a specific missing distinction.
Persistent Anomalies is the study of \emph{$\PhiF$ dynamics and
the geometry of the investigation frontier}.

\medskip
\noindent\textbf{Waves of Collapse.}
The Waves of Collapse program~\cite{flyxion2025d} studies the
discontinuous structure of archive compression: long plateaus of
normal encoding punctuated by collapse events.
In the present framework, this is the Devil's Staircase of
the total description length $\Lam$ under sequential template admissions.
Waves of Collapse is the study of \emph{collapse event dynamics and
the Devil's Staircase structure of $\Lam(t)$}.

\medskip
\noindent\textbf{Admissibility Geometry.}
The Admissibility Geometry program formalizes the admissibility
manifold and the geometry of projection failure.
It is the foundational mathematical framework from which the
present paper derives its formal content.
Admissibility Geometry is the study of \emph{the structure of
$\Adm_\tau$ and the metrics on the space of distinctions}.

\medskip
\noindent\textbf{Distinguishability Geometry.}
The Distinguishability Geometry program studies the metric structure
of the space of distinctions: how far apart are two cases, how
much information separates them, how the geometry changes under
different admissibility assumptions.
In the present framework, distinguishability geometry provides
the metric on the distinction space within which collapse events
are geodesic reorganizations.
Distinguishability Geometry is the study of \emph{the Riemannian
structure of the admissibility manifold}.

\medskip
\noindent\textbf{RSVP (Relativistic Scalar-Vector Plenum).}
The RSVP program provides the field-theoretic dynamics:
the equations governing how $\PhiF$, $v$ (rematching flow),
and $S$ (crystallized entropy) evolve, interact, and generate
shock fronts at collapse events.
RSVP is the study of \emph{the thermodynamics of distinction reorganization}.

\medskip
The six programs are cross-sections of the same geometric object:
the distinction structure of a domain under pressure from
accumulated residual and periodic reorganization by collapse events.
Each program captures a different aspect --- the obstruction to
collapse, the dynamics of anomaly accumulation, the structure of
discontinuous events, the geometry of admissibility, the metric
on distinctions, the thermodynamic field theory of reorganization ---
but each is addressing the same underlying phenomenon.

%% ============================================================
\section{Conclusion: The Master Claim and Its Corollaries}
\label{sec:conclusion}
%% ============================================================

The argument of this paper can be compressed into seven statements.

\begin{enumerate}
  \item \textbf{Projection destroys distinctions.}
        Every nontrivial representation collapses at least one
        distinction (Theorem~\ref{thm:projection}).
        Epistemology is the theory of which collapses matter.

  \item \textbf{Admissibility identifies which collapses matter.}
        A distinction is admissible for a task if collapsing it
        would impair navigation, intervention, or explanation.
        Projection failure is admissibility-relative:
        a harmless collapse for one task may be catastrophic for another.

  \item \textbf{Ontological deficit is invisible from within.}
        A distinction that is inexpressible in the current vocabulary
        manifests as structured residual $\PhiF$, not as an identifiable gap.
        The deficit becomes visible only from the successor vocabulary.

  \item \textbf{Residual density $\PhiF$ is the pressure toward reorganization.}
        Concentrated, structured $\PhiF$ defines the investigation frontier.
        Anomalies are not noise; they are the fuel for collapse events.

  \item \textbf{Understanding advances through collapse events.}
        A collapse event admits a new distinction $d^*$ that
        reorganizes the existing distinction structure, reduces
        $\bar\PhiF$ discontinuously, and enables rematching of the
        accumulated observational archive.
        Scientific revolutions, memory reconsolidation, archive
        rematching, and ontology repair in AI are all collapse events.

  \item \textbf{The World-Model Crossover and Directionality Principle
        mark the boundary between observation-centered and model-centered
        epistemology.}
        Below $\nstar$, observations organize theories.
        Above $\nstar$, theories organize observations.
        The byte becomes a unit of surprise rather than a unit of size.
        This is a principle and research program, not yet a theorem
        with a constructive proof in the general case.

  \item \textbf{Interpretability is distinction verification.}
        A model understands a domain to the degree that it preserves
        and makes recoverable the admissible distinctions.
        Chain-of-thought, mechanistic features, and internal
        representations are explanatory only insofar as they
        recover specific distinctions via recoverable operators.
\end{enumerate}

These seven statements converge on a single master claim:

\begin{quote}
\textbf{Understanding is not the accumulation of descriptions but
the reorganization of the distinctions from which descriptions
are constructed.}
\end{quote}

And because reorganization is irreducibly historical --- because
it occurs in time, requires a prior state to reorganize, and
produces a successor state that cannot be reached by accumulation
--- understanding is more precisely defined as a \emph{trajectory}:

\begin{quote}
\textbf{Understanding is a trajectory of distinction reorganization:
$\Dist_0 \to \Dist_1 \to \Dist_2 \to \cdots$,
in which each transition expands reachability, reduces residual density,
and retroactively remaps the accumulated observational archive.}
\end{quote}

This formulation resolves the tension between the static picture
(understanding as a property of a representation) and the dynamic
picture (understanding as a process). Both are correct at different
levels of description. The static picture characterizes a given
distinction structure $\Dist_t$ at a moment. The dynamic picture
characterizes the sequence and its properties. The static picture
is what a representation has; the dynamic picture is what an
understander does.

Prediction is not understanding because a predictor can collapse
admissible distinctions while achieving high accuracy.
Compression is not understanding because a compressor can reduce
size while discarding the distinctions required for intervention.
Description is not understanding because a description can be
accurate and inaccessible, faithful to the observation and
unfaithful to the domain.

Understanding is something more specific: a representation that
preserves the admissible distinctions of a domain, organizes them
into a structure with maximum reachability, and is capable of
reorganizing that structure when the residual field demands it.

The geometry of understanding is the geometry of this process.
It is not the geometry of facts in a space but the geometry of
distinctions in transformation.

\section*{Acknowledgements}

The RSVP field theory, Admissibility Geometry, Distinguishability
Geometry, Frozen Processes, and Persistent Anomalies programs
are developed across the author's ongoing independent research
corpus.
The Kuhnian translation in Section~\ref{sec:kuhn} benefited
from the observation that Kuhn's account becomes quantitative
once anomalies are identified with structured residual density.
The crossover theorem in Section~\ref{sec:crossover} is the
formal precipitate of ideas developed in the Waves of Collapse
and Rematching Archive programs.

\begin{thebibliography}{99}

\bibitem{bartlett1932}
Bartlett, F.~C. (1932).
\textit{Remembering: A Study in Experimental and Social Psychology}.
Cambridge University Press.

\bibitem{bengio2013}
Bengio, Y., Courville, A., \& Vincent, P. (2013).
Representation learning: a review and new perspectives.
\textit{IEEE Transactions on Pattern Analysis and Machine Intelligence},
35(8), 1798--1828.

\bibitem{friston2010}
Friston, K. (2010).
The free-energy principle: a unified brain theory?
\textit{Nature Reviews Neuroscience}, 11(2), 127--138.

\bibitem{jaynes2003}
Jaynes, E.~T. (2003).
\textit{Probability Theory: The Logic of Science}.
Cambridge University Press.

\bibitem{kolmogorov1965}
Kolmogorov, A.~N. (1965).
Three approaches to the quantitative definition of information.
\textit{Problems of Information Transmission}, 1(1), 1--7.

\bibitem{kuhn1962}
Kuhn, T.~S. (1962).
\textit{The Structure of Scientific Revolutions}.
University of Chicago Press.

\bibitem{lewis1973}
Lewis, D. (1973).
\textit{Counterfactuals}.
Harvard University Press.

\bibitem{loftus1979}
Loftus, E.~F. (1979).
\textit{Eyewitness Testimony}.
Harvard University Press.

\bibitem{nader2000}
Nader, K., Schafe, G.~E., \& LeDoux, J.~E. (2000).
Fear memories require protein synthesis in the amygdala for
reconsolidation after retrieval.
\textit{Nature}, 406, 722--726.

\bibitem{pearl2009}
Pearl, J. (2009).
\textit{Causality: Models, Reasoning, and Inference}.
Cambridge University Press.

\bibitem{rissanen1978}
Rissanen, J. (1978).
Modeling by shortest data description.
\textit{Automatica}, 14(5), 465--471.

\bibitem{shannon1948}
Shannon, C.~E. (1948).
A mathematical theory of communication.
\textit{Bell System Technical Journal}, 27, 379--423.

\bibitem{flyxion2025a}
Flyxion. (2025).
\textit{The Ecology of Distinctions}.
Independent monograph, 31 chapters.

\bibitem{flyxion2025b}
Flyxion. (2025).
\textit{Constraint, Projection, and Reachability}.
Independent textbook, 92 chapters.

\bibitem{flyxion2025c}
Flyxion. (2025).
\textit{Repair Pressure and the Lifecycle of Distinctions}.
Independent essay, 34 pages.

\bibitem{flyxion2025d}
Flyxion. (2025).
\textit{Waves of Collapse}.
Independent monograph.

\bibitem{flyxion2026a}
Flyxion. (2026).
\textit{Intelligence as Distinction Repair}.
Independent essay.

\bibitem{flyxion2026b}
Flyxion. (2026).
\textit{The Danger of Frozen Processes}.
Independent essay.

\bibitem{flyxion2026c}
Flyxion. (2026).
\textit{The Geometry of Distinction Preservation}.
Independent essay.

\end{thebibliography}

\end{document}
